\Subhead{2.5.1}{Уравнения эйконала и луча. Лучи и поле волны в слоисто-неоднородных средах}

Если показатель преломления меняется на масштабе $L\gg\lambda$, поле записывают $U=A(\bm r)\exp[\ii k_0S(\bm r)]$. Подстановка в $\Delta U+k_0^2n^2U=0$ и разделение порядков по $k_0L$ даёт
\begin{equation}
 |\nabla S|^2=n^2,\qquad
 2\nabla A\cdot\nabla S+A\Delta S=0,
 \quad\nabla\cdot(A^2\nabla S)=0.
\end{equation}
Первое — уравнение эйконала; второе выражает сохранение потока в лучевой трубке в скалярной недиспергирующей модели. Луч направлен по $\nabla S$; из принципа Ферма $\delta\int n\dd s=0$ получается
\begin{equation}\label{eq:rayequation}
 \frac{\dd}{\dd s}\left(n\frac{\dd\bm r}{\dd s}\right)=\nabla n.
\end{equation}
Поперечный градиент поворачивает луч в сторону большего $n$, то есть меньшей фазовой скорости. В поглощающей среде в уравнение потока добавляются потери.

При $n=n(z)$ горизонтальная компонента импульса сохраняется: $p=n\sin\theta=\mathrm{const}$, если $\theta$ — угол к вертикали. Поэтому
\begin{equation}
 \frac{\dd x}{\dd z}=\frac{p}{\sqrt{n^2(z)-p^2}},\qquad
 k_z=k_0\sqrt{n^2(z)-p^2}.
\end{equation}
Поворотная точка удовлетворяет $n(z_t)=p$. Вблизи неё простая лучевая амплитуда расходится; нужна волновая, например Airy-аппроксимация. На каустике пересекаются лучи, и геометрическая оптика также перестаёт быть равномерной. Она не описывает дифракционное заполнение тени без дополнительных построений.

У акустики роль $n$ играет величина, пропорциональная $1/c(z)$. Лучи изгибаются к области меньшей скорости. Это объясняет ионосферную рефракцию и звуковой канал, но конкретные профили сред определяют разные траектории.
